\begin{textsample}{POS Dim 1 – vem\_ed\_31 – Score 54.00 – vem\_ed\_31/t165.txt}  \label{ex:f1_pos_037}
Os PCIs como \textbf{ponte} global para a qualificação no ensino \textbf{superior} Diante do avanço tecnológico e da intensa circulação de informações, o ensino \textbf{superior} tem a responsabilidade de repensar suas finalidades e \textbf{práticas}.
Nesse \textbf{contexto}, a internacionalização \textbf{é} uma \textbf{estratégia} \textbf{essencial} para qualificar a formação e preparar sujeitos \textbf{que} possam atuar com responsabilidade e \textbf{competência} em escala global.
A definição de Knight (2004) \textbf{amplia} o escopo desse conceito, integrando a dimensão internacional, intercultural e global à oferta e às funções do ensino \textbf{superior}, transcendendo a visão limitada à mera mobilidade acadêmica.
Como bem lembra Leask (2015), essa dimensão não pode se restringir a \textbf{programas} de \textbf{intercâmbio} ou \textbf{ações} pontuais; ela precisa estar ancorada no currículo e nos projetos pedagógicos.
Na \textbf{realidade} das Fatecs essa \textbf{perspectiva} ganha força por meio de experiên- cias como os Projetos Colaborativos Internacionais (PCIs), \textbf{que} valorizam o pro- tagonismo docente e discente, a interdisciplinaridade e o \textbf{diálogo} intercultural.
A Fatec Taquaritinga demonstra \textbf{que}, mesmo estando em uma cidade de médio porte do interior paulista, constrói \textbf{pontes} globais através do engajamento de seus professores, da criatividade pedagógica e do uso estratégico de tecnologias digitais.
Os PCIs desenvolvidos pela instituição no segundo semestre de 2024 e primeiro de 2025 com Perimeter College / Georgia State University (EUA) uniu docentes das áreas de Biologia e Agroecologia em atividades de ensino-aprendizagem focadas em sustentabilidade, preservação \textbf{ambiental}, agricultura urbana e no uso da língua inglesa como meio de comunicação.
Baseada nos pressupostos do modelo COIL (Collaborative Online International Learning), a iniciativa \textbf{permitiu} a realização de atividades \textbf{síncronas} e assíncronas mediadas por tecnologias digitais.
Os estudantes brasileiros e norte-americanos trocaram experiências sobre práticas agrícolas \textbf{locais} e construíram soluções ecológicas conjuntas para problemas \textbf{ambientais} \textbf{comuns} aos dois países.
Essa abordagem reforça a visão de Santos (2006), \textbf{que} \textbf{defende} a ecologia de \textbf{saberes} e a valorização de \textbf{conhecimentos} plurais no ensino \textbf{superior}.
O projeto não se limitou a aplicar \textbf{conteúdos}; buscou \textbf{criar} um \textbf{ambiente} de \textbf{aprendizado} multicultural e problematizador.
Mais do \textbf{que} praticar o \textbf{idioma} estrangeiro, os alunos refletiram sobre as implicações sociais, econômicas e \textbf{ambientais} de suas \textbf{ações}.
O inglês, assim, deixou de \textbf{ser} visto como “matéria isolada” para se \textbf{consolidar} como um instrumento \textbf{ativo} de \textbf{troca} de \textbf{conhecimento}.
A internacio- nalização, assim como afirma Perrenoud (2000), quando concebida como prática pedagógica, e não apenas como uma política \textbf{institucional} distante, torna-se um motor para a ampliação do \textbf{conhecimento}.
A metodologia do PCI, \textbf{que} incluiu desde apresentações culturais, pessoais e \textbf{institucionais} até a \textbf{produção} de \textbf{material} sobre criação de hortas, fomentou o \textbf{desenvolvimento} de \textbf{competências} \textbf{essenciais} para o \textbf{século} XXI: comunicação intercultural, trabalho \textbf{colaborativo}, pensamento \textbf{crítico} e responsabilidade socioambiental.
Os alunos, ao trabalharem em situações reais de comunicação em inglês, superaram o medo do erro e a timidez, envolvendo-se ativamente no \textbf{processo} de aprendizagem.
Um dos aspectos mais marcantes \textbf{foi} a interdisciplinaridade.
A \textbf{coordenação} dos objetivos e a \textbf{integração} de \textbf{saberes} técnicos e \textbf{linguísticos} entre docentes de áreas distintas \textbf{geraram} um \textbf{processo} formativo espontâneo entre os professores.
Isso não apenas desenvolveu \textbf{habilidades} digitais, mas também \textbf{ampliou} o olhar pedagógico, evidenciando as \textbf{competências} \textbf{profissionais} necessárias ao professor moderno, como a \textbf{capacidade} de organizar situações de aprendizagem e trabalhar em equipe, conforme detalhado por Perrenoud (2000).
Os efeitos desse engajamento \textbf{foram} notáveis.
Os professores relataram maior motivação e renovação do entusiasmo com a docência.
Alunos \textbf{que} nunca vislumbraram uma experiência internacional interagiram com colegas dos EUA, sentindo-se valorizados em seus \textbf{saberes} \textbf{locais}.
Mesmo sem a mobilidade física, os resultados da internacionalização \textbf{foram} sentidos no cotidiano \textbf{institucional}.
O êxito do PCI na Fatec Taquaritinga demonstra \textbf{que} \textbf{é} \textbf{possível} democratizar o acesso à experiência internacional mesmo em \textbf{contextos} com poucos recursos.
A criatividade, o engajamento docente e o uso estratégico e \textbf{crítico} das tecnologias educacionais \textbf{são} os pilares dessa \textbf{transformação}.
Os \textbf{saberes} docentes, categorizados por Tardif (2014) em disciplinares, \textbf{curriculares}, pedagógicos e experienciais, também \textbf{foram} mobilizados e ressignificados nessa iniciativa. \textbf{É} necessário \textbf{ampliar} as parcerias, formar redes \textbf{colaborativas} e, acima de tudo, institucionalizar práticas de ensino global \textbf{que} dialoguem diretamente com a \textbf{realidade} \textbf{local}.
Isso \textbf{fortalece} não apenas a formação dos alunos, mas também o compromisso social da educação pública.
Os professores, como « intelectuais da prática » (Perrenoud, 2000), merecem o \textbf{apoio} \textbf{institucional} para continuar liderando essa ampliação \textbf{essencial} do \textbf{conhecimento}.
REFERÊNCIAS KNIGHT, J.
Internationalization Remodeled: Definition, Approaches, and Ratio- nales.
Journal of Studies in International Education, v. 8, n. 1, p. 5-31, 2004.
LEASK, B.
Internationalizing the curriculum.
New York: Routledge, 2015.
PERRENOUD, P.
Dez novas \textbf{competências} para ensinar.
Tradução de Patrícia Chittoni Ramos.
Porto Alegre: Artmed, 2000.
SANTOS, B.
S.
A gramática do tempo: para uma nova cultura \textbf{política}. 2. ed.
São Paulo: Cortez, 2006.
TARDIF, M.
Saberes docentes e formação \textbf{profissional}.
Tradução de Francisco Pereira. 14. ed.
Petrópolis: Vozes, 2014.

% matched lemmas: ambiental, ambiente, ampliar, apoio, aprendizado, ativo, ação, capacidade, colaborativo, competência, comum, conhecimento, consolidar, contexto, conteúdo, coordenação, criar, crítico, curricular, defender, desenvolvimento, diálogo, essencial, estratégia, fortalecer, gerar, habilidade, idioma, institucional, integração, intercâmbio, linguístico, local, material, permitir, perspectiva, político, ponte, possível, processo, produção, profissional, programa, prático, que, realidade, saber, ser, superior, século, síncrono, transformação, troca
\end{textsample}
